% Chapter 3, Figure 32: drag coefficient of a paraboloidal nose against its fineness ratio L/d (scan:
% figures/ch3/fig32.png; after Stine). The curve is digitized from the scan by fig32.py (no equation in the
% book); overlay on the scan (fig32.calib.json): 95% of the points within 0 px of the ink, 95% of the traced
% ink centre line within 0.39 px of the curve.
% The inset is drawn in the scan's pixels at the plot's scale (the axes are 443.5 x 288 px on the scan, 4.2 in
% wide here): a paraboloid of revolution (r ~ sqrt(x) from the tip, as the scan's outline: L = 219 px,
% d = 55 px, L/d = 4) with its flat base, on its centre line, L dimensioned below and d at the right.
\documentclass[11pt]{standalone}
\usepackage{tamrfig}
\begin{document}
\begin{tikzpicture}
  \def\u{0.024054cm}   % one scan pixel
  \begin{axis}[tamr,
      width=443.5*\u, height=288*\u,
      xmin=0, xmax=4, xtick={0,...,4}, minor xtick={0.5,1.5,2.5,3.5},
      ymin=0, ymax=2, ytick={0,0.5,...,2}, yticklabels={0,0.5,1.0,1.5,2.0},
      minor ytick={0.25,0.75,1.25,1.75},
      xlabel={$\dfrac{L}{d}$},
      ylabel={$\CD$}, ylabel style={rotate=-90, anchor=east}]
    \addplot[series1] table[x=x, y=y, col sep=comma] {figures/v2/ch3/fig32.csv};
    \coordinate (O) at (axis cs:0,0);
  \end{axis}
  % the inset, in the scan's pixels (y down) from its axes' origin (88, 305.5)
  \begin{scope}[shift={(O)}, x=\u, y=-\u, xshift=-88*\u, yshift=305.5*\u]
    \draw[centerline] (262,124.5) -- (508,124.5);
    \begin{scope}[shift={(276,124.5)}]
      \rocketoutline[parabolic]{219/27.5}
    \end{scope}
    % L below
    \draw[extension] (276,129) -- (276,186);
    \draw[extension] (495,155) -- (495,186);
    \dimline{(276,179)}{(495,179)}{$L$}
    % d at the right
    \draw[extension] (499,97) -- (528,97);
    \draw[extension] (499,152) -- (528,152);
    \dimline{(520,97)}{(520,152)}{$d$}
  \end{scope}
\end{tikzpicture}
\end{document}
